\DocumentMetadata{
  pdfversion=2.0,
  pdfstandard=ua-2,
  lang=en-US,
  tagging=on,
  tagging-setup={math/setup=mathml-SE,tabsorder=structure}
}
\documentclass[11pt,letterpaper]{article}
\usepackage[margin=0.68in,includefoot]{geometry}
\usepackage{newcomputermodern}
\usepackage{amsmath,amscd,mathtools}
\usepackage{tikz,adjustbox}
\providecommand{\square}{\mdlgwhtsquare}
\usepackage[hidelinks]{hyperref}
\hypersetup{
  pdftitle={Algebra PhD Exam, May 2014},
  pdfauthor={University of Florida Department of Mathematics},
  pdfsubject={Graduate mathematics examination},
  pdfdisplaydoctitle=true,
  bookmarksopen=true
}
\setcounter{secnumdepth}{0}
\setlength{\parindent}{0pt}
\setlength{\parskip}{0.28em}
\setlength{\emergencystretch}{3em}
\setlength{\leftmargini}{2.15em}
\setlength{\leftmarginii}{2.65em}
\setlength{\leftmarginiii}{3em}
\pagestyle{empty}
\raggedbottom
\allowdisplaybreaks
\NewTaggingSocketPlug{block/list/label}{examproblem}{%
  \tagstructbegin{tag=Lbl}\tagstructend%
  \tagstructbegin{tag=\LBody}%
  \tagstructbegin{tag=\ProblemHeadingTag,title-o={\ProblemHeadingText}}%
  \tagmcbegin{tag=\ProblemHeadingTag,actualtext={\ProblemHeadingText}}%
  #2%
  \tagmcend%
  \tagstructend%
}
\newenvironment{examproblems}[2]{%
  \def\ProblemHeadingTag{#2}%
  \AssignTaggingSocketPlug{block/list/label}{examproblem}%
  \begin{enumerate}%
  \setcounter{enumi}{#1}%
  \renewcommand{\labelenumi}{\textbf{\theenumi.}}%
  \setlength{\topsep}{0.35em}%
  \setlength{\partopsep}{0pt}%
  \setlength{\itemsep}{0.48em}%
  \setlength{\parsep}{0.18em}%
}{\end{enumerate}}
\newenvironment{examsubparts}{%
  \AssignTaggingSocketPlug{block/list/label}{default}%
  \begin{enumerate}%
  \setlength{\topsep}{0.18em}%
  \setlength{\partopsep}{0pt}%
  \setlength{\itemsep}{0.16em}%
  \setlength{\parsep}{0.1em}%
}{\end{enumerate}}
\newenvironment{examinstructions}{%
  \par\begingroup\itshape
}{%
  \par\endgroup\vspace{0.18em}
}
\makeatletter
\renewcommand\subsection{\@startsection{subsection}{2}{0pt}%
  {-1.25ex plus -.25ex minus -.1ex}%
  {0.55ex plus .1ex}%
  {\normalfont\large\bfseries}}
\renewcommand{\@maketitle}{%
  \newpage\null\vskip -1em%
  {\footnotesize\bfseries
    \href{https://www.ufl.edu/}{University of Florida}, \href{https://math.ufl.edu/}{Department of Mathematics}\par}%
  \vskip -0.5em%
  \tagstructbegin{tag=H1,title={Algebra PhD Exam, May 2014}}%
  \tagpdfparaOff%
  \tagmcbegin{tag=H1}%
    {\LARGE\bfseries\href{https://gma.math.ufl.edu/exams/algebra-phd/}{Algebra PhD Exam}, May 2014\par}%
  \tagmcend%
  \tagpdfparaOn%
  \tagstructend%
  \vskip 0.25em%
  \tagmcbegin{artifact}%
    \hrule height 1pt%
  \tagmcend%
  \vskip 1em%
}
\makeatother
\title{Algebra PhD Exam, May 2014}
\author{}
\date{}
\begin{document}
\maketitle
\thispagestyle{empty}
\begin{examinstructions}
Do any 7 of the 11 problems, and mark clearly which problems you want graded if you try more than 7. State clearly any theorems that you use.
\end{examinstructions}
\begin{examproblems}{0}{H2}
\def\ProblemHeadingText{Problem 1.}
\item
Suppose \(p_1,p_2,\ldots,p_r\) are distinct primes, and \(K=\mathbb{Q}(\sqrt{p_1},\ldots,\sqrt{p_r})\), find the Galois group of \(K/\mathbb{Q}\). How many fields~\(F\) are there such that \(\mathbb{Q}\subseteq F\subseteq K\)?
\def\ProblemHeadingText{Problem 2.}
\item
This problem has two parts.
\begin{examsubparts}
\item[\textnormal{(i)}]
Show that if \(K/k\) and \(L/K\) are separable algebraic extensions, then so is \(L/k\).
\item[\textnormal{(ii)}]
Give an example of normal extensions \(K/k\), \(L/K\) such that \(L/k\) is not normal.
\end{examsubparts}
\def\ProblemHeadingText{Problem 3.}
\item
Suppose~\(K\) is an algebraic extension field of~\(k\) and \(\bar{k}\) is an algebraic closure of~\(k\). Show that \(K/k\) is separable if and only if the ring \(\bar{k}\otimes_k K\) has no nonzero nilpotent elements.
\def\ProblemHeadingText{Problem 4.}
\item
Suppose~\(A\) is a commutative ring with identity and \(S\subset A\) is a multiplicative set not containing~\(0\).
\begin{examsubparts}
\item[\textnormal{(i)}]
Show that there is a prime ideal \(\mathfrak{p}\subset A\) such that \(\mathfrak{p}\cap S=\phi\).
\item[\textnormal{(ii)}]
Use part (i) to show that the intersection of all prime ideals of~\(A\) is the set of nilpotent elements.
\end{examsubparts}
\def\ProblemHeadingText{Problem 5.}
\item
Suppose~\(k\) is a field and~\(A\) is a finitely generated commutative~\(k\)-algebra. Let \(\bar{k}\) be an algebraic closure of~\(k\). Note that the Galois group \(G=\operatorname{Gal}(\bar{k}/k)\) acts on the set \(\operatorname{Hom}_k(A,\bar{k})\) of~\(k\)-algebra homomorphisms \(A\to\bar{k}\) by composition. Construct a bijection of the set of orbits of~\(G\) on \(\operatorname{Hom}_k(A,\bar{k})\) with the set of maximal ideals in~\(A\).
\def\ProblemHeadingText{Problem 6.}
\item
Suppose~\(A\) is an integral domain. Show that an ideal \(I\subseteq A\) is invertible if and only if it is projective and finitely generated.
\def\ProblemHeadingText{Problem 7.}
\item
Suppose~\(X\) is a nonempty set and \(F(X)\) is the free group on~\(X\). Show that \(F(X)\) is abelian if and only if~\(X\) is a singleton.
\def\ProblemHeadingText{Problem 8.}
\item
This problem has two parts.
\begin{examsubparts}
\item[\textnormal{(i)}]
Define what it means for a group to be nilpotent.
\item[\textnormal{(ii)}]
Show that if~\(G\) is a finite nilpotent group, then any Sylow subgroup is normal.
\end{examsubparts}
\def\ProblemHeadingText{Problem 9.}
\item
Suppose~\(A\) is a commutative ring with identity and~\(B\) is a commutative~\(A\)-algebra with structure map \(f:A\to B\). Show that the forgetful functor from the category of~\(B\)-modules to the category of~\(A\)-modules induced by~\(f\) has a left adjoint.
\def\ProblemHeadingText{Problem 10.}
\item
Let~\(p\) be a prime. Find all the isomorphism classes of noncommutative semisimple rings~\(A\) of cardinality \(p^{10}\).
\def\ProblemHeadingText{Problem 11.}
\item
Suppose~\(R\) is a finite commutative ring with identity, such that \(x^3=x\) for all \(x\in R\). Show that~\(R\) is a finite direct sum of fields isomorphic to \(\mathbb{F}_2\) or \(\mathbb{F}_3\).
\end{examproblems}
\end{document}
